% Generated draft: requires mathematical and source review.
\documentclass[11pt]{article}
\usepackage[margin=25mm]{geometry}
\usepackage{fontspec}
\setmainfont{Latin Modern Roman}
\usepackage{amsmath,amssymb,mathtools,unicode-math}
\setmathfont{Latin Modern Math}
\usepackage{enumitem,needspace,xurl,hyperref}
\setlength{\parindent}{0pt}
\setlength{\parskip}{5pt}
\setlength{\emergencystretch}{3em}
\newcommand{\N}{\mathbb{N}}
\newcommand{\Z}{\mathbb{Z}}
\newcommand{\Q}{\mathbb{Q}}
\newcommand{\R}{\mathbb{R}}
\newcommand{\C}{\mathbb{C}}
\newcommand{\F}{\mathbb{F}}
\newcommand{\A}{\mathbb{A}}
\newcommand{\PP}{\mathbb{P}}
\newcommand{\E}{\mathbb{E}}
\newcommand{\eps}{\varepsilon}
\newcommand{\dd}{\,\mathrm{d}}
\begin{document}
\section*{UnsolvedMath ID: 30005953}
\textbf{Source problem code:} catalogue entry
\subsection{Definitions and mathematical statement}
Let \(S^4=\{x\in\R^5:\|x\|=1\}\) be the standard smooth four-sphere. Write \(\pi_0\operatorname{Diff}(S^4)\) for the group of smooth isotopy classes of smooth diffeomorphisms, with multiplication induced by composition and identity \(1=[\operatorname{id}_{S^4}]\). The supplied excerpt does not specify whether an orientation-preserving restriction is intended; that convention requires confirmation from the report.

The supplied record identifies a particular class by the reported equalities
\[
\alpha:=\operatorname{ps}\circ r([c])
=\operatorname{wat}(\Theta)
=\operatorname{impl}(yx)
=G(\nu T_c)
\in\pi_0\operatorname{Diff}(S^4),
\qquad \alpha^2=1.
\]
Here ``at-most-two-torsion'' means that the order of \(\alpha\) divides two; it does not assert that \(\alpha\) is nontrivial. The question is whether
\[
\alpha=1,
\]
that is, whether a representative of this particular class is smoothly isotopic to the identity.

The symbols \(\operatorname{ps},r,[c],\operatorname{wat},\Theta,\operatorname{impl},y,x,G,\nu T_c\) are retained as source notation. Their constructions, domains, auxiliary choices, and the hypotheses of the referenced theorem are absent from the supplied excerpt. Thus these expressions identify the intended class only by reference to that theorem; they do not give a self-contained definition. The original excerpt describes Figure 3 as producing Watanabe's \(\operatorname{wat}(\Theta)\) and Budney--Gabai's \(\operatorname{impl}(yx)\) by parameterised surgery on depicted families of circles. Neither those families nor the figure are supplied.
\subsection{Short English statement}
The question asks whether the particular parametrized-surgery class on \(S^4\) identified by the four reported expressions is smoothly isotopic to the identity. Its reported order divides two. The constructions defining the class and the referenced theorem must be recovered from the report before this draft is mathematically self-contained.
\subsection{Sources}
\begin{itemize}
\item Supplied local UnsolvedMath registry record, permanent ID \(30005953\), source problem code \texttt{catalogue entry}, entitled ``Triviality of the Parametrized-Surgery Class on \(S^4\).'' No proposer is named; the supplied proposed year is 2024.
\item Supplied report citation: \emph{Topologie} (2025), \emph{Oberwolfach Reports}, pp.~1913--1974, DOI \texttt{10.4171/owr/2024/34}. Source URL and supplied literature reference labelled ``original OWR report'': \url{https://doi.org/10.4171/owr/2024/34}.
\item The original excerpt attributes \(\operatorname{wat}(\Theta)\) to Watanabe and \(\operatorname{impl}(yx)\) to Budney--Gabai, and refers to a theorem and Figure 3 without supplying them. No separate publications for these constructions are supplied.
\item Literature assessment dated exactly \texttt{2026-08-22}: the supplied curation labels the problem open, says that the report explicitly leaves triviality unresolved, and reports that no later primary resolution was located. This is supplied dated curation, not an independently verified literature assessment. The record also labels its statement \texttt{corrected\_verified} and reports verification against the official report; that verification claim has not been independently checked here.
\end{itemize}
\end{document}
